\SourceLeaf{8}
\begin{center}
\textsc{DEM DURCHLAUCHTIGSTEN}

\smallskip

\textsc{FÜRSTEN UND HERRN}

\smallskip

{\Large\textsc{KARL WILHELM FERDINAND}}

\smallskip

\textsc{HERZOG VON BRAUNSCHWEIG UND LÜNEBURG.}

\smallskip

\BookRule{0.25\linewidth}

\smallskip

\textsc{DURCHLAUCHTIGSTER FÜRST}
\end{center}

\noindent Ich schätze mich sehr glücklich, dass \textsc{Ihr} mir erlaubt, dieses Werk \textsc{Eurem} erhabenen Namen zu widmen; eine heilige Pflicht der Dankbarkeit verbindet mich, es \textsc{Euch} darzubringen. Hätte mir nicht \textsc{Eure} Gunst, durchlauchtigster Fürst, zuerst den Zugang zu den Wissenschaften eröffnet, hätten nicht \textsc{Eure} beständigen Wohltaten meine Studien immer unterstützt, so hätte ich mich nicht ganz der Mathematik widmen können, zu der mich stets eine leidenschaftliche Liebe gezogen hat. Ja, dass ich gerade die Untersuchungen beginnen, über mehrere Jahre fortsetzen und niederschreiben konnte, von denen dieser Band einen Teil enthält, hat allein \textsc{Eure} Güte ermöglicht: Sie befreite mich von anderen Sorgen, sodass ich mich vor allem dieser Arbeit widmen konnte. Als ich sie schließlich veröffentlichen wollte, beseitigte \textsc{Eure} Freigebigkeit alle Hindernisse, die den Druck verzögerten. Diese großen Verdienste, die \textsc{Ihr} Euch um mich und meine Bemühungen erworben habt, kann ich eher mit dankbarem Herzen und stiller Bewunderung \SourceLeaf{9}bedenken, als sie mit gerechtem und würdigem Lob feiern. Denn ich fühle mich nicht nur einer solchen Aufgabe nicht gewachsen, sondern glaube auch, dass jeder weiß, wie gewohnt \textsc{Ihr} seid, allen, die zur Pflege der besten Wissenschaften beizutragen scheinen, eine so außergewöhnliche Freigebigkeit zu erweisen. Auch die Wissenschaften, die man gewöhnlich für schwerer zugänglich und vom Nutzen des täglichen Lebens weiter entfernt hält, sind von \textsc{Eurem} Schutz nicht ausgeschlossen: \textsc{Ihr selbst} habt mit jenem weisen Geist, der alles so gut kennt, was zum Wohlergehen der menschlichen Gesellschaft beiträgt, den inneren und notwendigen Zusammenhang aller Wissenschaften vollständig erkannt. Wenn \textsc{Ihr}, durchlauchtigster Fürst, dieses Buch, das sowohl von meiner tiefen Dankbarkeit gegen \textsc{Euch} als auch von meiner einer höchst edlen Wissenschaft gewidmeten Arbeit zeugt, jener besonderen Gunst nicht unwürdig findet, mit der Ihr mich so lange begleitet habt, werde ich mich glücklich schätzen, meine Mühe nicht vergebens aufgewandt und die Ehre erlangt zu haben, die ich mir vor allen anderen gewünscht habe

\begin{center}
\textsc{DURCHLAUCHTIGSTER FÜRST}
\end{center}

\noindent\hspace*{1.5em}Braunschweig, im Juli 1801.

\smallskip
\begin{flushright}
Eurer Durchlaucht ergebenster Diener

\smallskip
\textsc{C. F. Gauss.}
\end{flushright}
